\section{文本宏}

\subsection{强调命令 (dfntxt)}

\dfntxt{特征值分解} 是将矩阵分解为特征值和特征向量的过程。

\dfntxt{奇异值分解 (SVD)} 是特征值分解的推广，可以应用于任意矩阵。

\dfntxt{主成分分析 (PCA)} 是一种常用的数据降维方法，基于奇异值分解。

\dfntxt{条件数} 衡量了矩阵对扰动的敏感程度。条件数越大，矩阵越病态。

\subsection{列表项缩进 (tabitem)}

\tabitem 向量空间：满足八条公理的集合
\tabitem 线性变换：保持加法和数乘的映射
\tabitem 矩阵：线性变换的表示
\tabitem 特征值：描述线性变换的伸缩因子
\tabitem 特征向量：在变换下方向不变的向量

\tabitem 监督学习：使用标注数据训练模型
\tabitem 无监督学习：发现数据的内在结构
\tabitem 强化学习：通过试错学习最优策略
\tabitem 半监督学习：结合标注和未标注数据